\documentclass[12pt,a4paper]{report}

% ============================================
% Paquetes
% ============================================
\usepackage[utf8]{inputenc}
\usepackage[T1]{fontenc}
\usepackage[spanish]{babel}
\usepackage{geometry}
\usepackage{graphicx}
\usepackage{xcolor}
\usepackage{listings}
\usepackage{hyperref}
\usepackage{booktabs}
\usepackage{tabularx}
\usepackage{float}
\usepackage{amsmath}
\usepackage{enumitem}
\usepackage{fancyhdr}
\usepackage{titlesec}
\usepackage{caption}
\usepackage{subcaption}
\usepackage{tikz}
\usetikzlibrary{shapes,arrows.meta,positioning,fit,backgrounds}

\geometry{margin=2.5cm}

% ============================================
% Colores
% ============================================
\definecolor{codegreen}{rgb}{0.0, 0.5, 0.0}
\definecolor{codegray}{rgb}{0.5, 0.5, 0.5}
\definecolor{codepurple}{rgb}{0.58, 0.0, 0.82}
\definecolor{backcolour}{rgb}{0.97, 0.97, 0.97}
\definecolor{dbgreen}{HTML}{22c55e}
\definecolor{dbyellow}{HTML}{eab308}
\definecolor{dbred}{HTML}{ef4444}
\definecolor{linkblue}{HTML}{2563eb}

% ============================================
% Configuracion de listings
% ============================================
\lstdefinestyle{code}{
    backgroundcolor=\color{backcolour},
    commentstyle=\color{codegreen},
    keywordstyle=\color{codepurple}\bfseries,
    numberstyle=\tiny\color{codegray},
    stringstyle=\color{red!70!black},
    basicstyle=\ttfamily\footnotesize,
    breakatwhitespace=false,
    breaklines=true,
    captionpos=b,
    keepspaces=true,
    numbers=left,
    numbersep=8pt,
    showspaces=false,
    showstringspaces=false,
    showtabs=false,
    tabsize=2,
    frame=single,
    rulecolor=\color{gray!30},
    xleftmargin=1.5em,
    framexleftmargin=1.5em,
}

\lstdefinestyle{json}{
    style=code,
    morestring=[b]",
    literate=
     *{0}{{{\color{codepurple}0}}}{1}
      {1}{{{\color{codepurple}1}}}{1}
      {2}{{{\color{codepurple}2}}}{1}
      {3}{{{\color{codepurple}3}}}{1}
      {4}{{{\color{codepurple}4}}}{1}
      {5}{{{\color{codepurple}5}}}{1}
      {6}{{{\color{codepurple}6}}}{1}
      {7}{{{\color{codepurple}7}}}{1}
      {8}{{{\color{codepurple}8}}}{1}
      {9}{{{\color{codepurple}9}}}{1}
      {:}{{{\color{codegray}:}}}{1}
      {,}{{{\color{codegray},}}}{1},
}

\lstset{style=code}

% ============================================
% Cabeceras
% ============================================
\pagestyle{fancy}
\fancyhf{}
\fancyhead[L]{\leftmark}
\fancyhead[R]{Monitor de Ruido en Aulas}
\fancyfoot[C]{\thepage}
\renewcommand{\headrulewidth}{0.4pt}

% ============================================
% Hyperref
% ============================================
\hypersetup{
    colorlinks=true,
    linkcolor=linkblue,
    urlcolor=linkblue,
    citecolor=linkblue,
    pdfauthor={IES Enlaces},
    pdftitle={Monitor de Ruido en Aulas - Documentacion Completa},
}

% ============================================
% Titulo
% ============================================
\title{
    \vspace{-2cm}
    {\Huge\bfseries Monitor de Ruido en Aulas}\\[0.5cm]
    {\Large Sistema de monitorizaci\'on ac\'ustica en tiempo real}\\[0.5cm]
    {\large Documentaci\'on t\'ecnica completa}\\[1cm]
    \rule{\textwidth}{1pt}\\[0.5cm]
    {\normalsize
    \begin{tabular}{ll}
    \textbf{Centro:} & IES Enlaces \\
    \textbf{Repositorio:} & \url{https://github.com/Rubenbros/microfonoAula} \\
    \textbf{Stack:} & ESP32 + MQTT + Node.js + Next.js \\
    \textbf{Fecha:} & Marzo 2026 \\
    \end{tabular}
    }
}
\author{}
\date{}

\begin{document}

\maketitle
\thispagestyle{empty}

\tableofcontents
\newpage

% ============================================================
% CAPITULO 1: Introduccion
% ============================================================
\chapter{Introducci\'on}

\section{Descripci\'on del proyecto}

El \textbf{Monitor de Ruido en Aulas} es un sistema de monitorizaci\'on ac\'ustica en tiempo real dise\~nado para medir, registrar y visualizar los niveles de ruido en aulas escolares. El sistema utiliza microcontroladores M5Stack (basados en ESP32) equipados con micr\'ofonos MEMS integrados que capturan audio, aplican filtrado A-weighting conforme a la norma IEC~61672, y transmiten las mediciones v\'ia MQTT a un servidor central.

El objetivo principal es proporcionar al profesorado y equipo directivo una herramienta para:

\begin{itemize}
    \item Monitorizar en tiempo real el nivel de ruido de cada aula.
    \item Detectar picos de ruido y patrones temporales.
    \item Comparar niveles entre diferentes zonas del aula (micr\'ofonos distribuidos).
    \item Almacenar hist\'orico de mediciones para an\'alisis posterior.
    \item Visualizar los datos en un dashboard web accesible desde cualquier dispositivo.
\end{itemize}

\section{Arquitectura general}

El sistema sigue una arquitectura IoT cl\'asica de tres capas:

\begin{figure}[H]
\centering
\begin{tikzpicture}[
    node distance=1.5cm and 2.5cm,
    box/.style={draw, rounded corners, minimum height=1.2cm, minimum width=3cm, align=center, fill=#1!15, font=\small},
    arrow/.style={-{Stealth[length=3mm]}, thick},
]
    % Capa dispositivos
    \node[box=blue] (atom1) {ATOM Echo\\mic\_01..06};
    \node[box=blue, right=2cm of atom1] (core2) {Core2\\mic\_central};

    % Broker
    \node[box=orange, below=2cm of $(atom1)!0.5!(core2)$] (mqtt) {Broker MQTT\\Puerto 1883};

    % Backend
    \node[box=green, below=2cm of mqtt] (backend) {Backend Node.js\\Express + SQLite};

    % Frontend
    \node[box=purple, below=2cm of backend] (frontend) {Frontend Next.js\\Dashboard React};

    % Flechas
    \draw[arrow] (atom1) -- node[left, font=\scriptsize] {WiFi} (mqtt);
    \draw[arrow] (core2) -- node[right, font=\scriptsize] {WiFi} (mqtt);
    \draw[arrow, <->] (mqtt) -- node[right, font=\scriptsize] {MQTT} (backend);
    \draw[arrow, <->] (backend) -- node[right, font=\scriptsize] {REST + WS} (frontend);

    % Labels
    \node[left=0.5cm of atom1, font=\footnotesize\itshape, text=gray] {Capa f\'isica};
    \node[left=0.5cm of mqtt, font=\footnotesize\itshape, text=gray] {Mensajer\'ia};
    \node[left=0.5cm of backend, font=\footnotesize\itshape, text=gray] {Servidor};
    \node[left=0.5cm of frontend, font=\footnotesize\itshape, text=gray] {Visualizaci\'on};
\end{tikzpicture}
\caption{Arquitectura general del sistema}
\end{figure}

\section{Stack tecnol\'ogico}

\begin{table}[H]
\centering
\begin{tabularx}{\textwidth}{lXl}
\toprule
\textbf{Capa} & \textbf{Tecnolog\'ia} & \textbf{Versi\'on} \\
\midrule
Firmware & C++ (Arduino) + PlatformIO & ESP-IDF 5.x \\
Comunicaci\'on & MQTT (Aedes / Mosquitto) & 5.x \\
Backend & Node.js + Express + SQLite & 20 LTS \\
Frontend & Next.js + React + TypeScript & 15.1 / 19.0 \\
Visualizaci\'on & Recharts + Tailwind CSS & 2.15 / 3.4 \\
Base de datos & SQLite (better-sqlite3) & 11.7 \\
\bottomrule
\end{tabularx}
\caption{Stack tecnol\'ogico del proyecto}
\end{table}

% ============================================================
% CAPITULO 2: Hardware
% ============================================================
\chapter{Hardware}

\section{M5Stack ATOM Echo S3R (Micr\'ofono distribuido)}

El ATOM Echo S3R es el dispositivo principal del sistema. Se despliegan 6 unidades por aula en posiciones estrat\'egicas para cubrir diferentes zonas.

\begin{table}[H]
\centering
\begin{tabularx}{\textwidth}{lX}
\toprule
\textbf{Caracter\'istica} & \textbf{Especificaci\'on} \\
\midrule
MCU & ESP32-S3 (240 MHz dual-core, PSRAM) \\
Micr\'ofono & PDM MEMS integrado \\
Interfaz audio & I2S \\
Conectividad & WiFi 802.11 b/g/n (2.4 GHz) \\
LED & RGB integrado (indicador de estado) \\
Bot\'on & Bot\'on A (calibraci\'on / reset pico) \\
Alimentaci\'on & USB-C (5V) \\
Tama\~no & 24 $\times$ 24 $\times$ 17 mm \\
\bottomrule
\end{tabularx}
\caption{Especificaciones del ATOM Echo S3R}
\end{table}

\subsection{Configuraci\'on GPIO (I2S)}

\begin{table}[H]
\centering
\begin{tabular}{lll}
\toprule
\textbf{Se\~nal} & \textbf{GPIO} & \textbf{Funci\'on} \\
\midrule
LRCK & GPIO 3 & Left/Right Clock \\
SCLK & GPIO 17 & Serial Clock \\
MCLK & GPIO 11 & Master Clock \\
DSDIN & GPIO 48 & Data Serial In \\
\bottomrule
\end{tabular}
\caption{Pines I2S del ATOM Echo S3R}
\end{table}

\subsection{Indicador LED}

El LED RGB indica el nivel de ruido actual:

\begin{itemize}
    \item \textcolor{dbgreen}{\textbf{Verde}}: $< 50$ dBA (ambiente tranquilo)
    \item \textcolor{dbyellow}{\textbf{Amarillo}}: $50-70$ dBA (nivel moderado)
    \item \textcolor{dbred}{\textbf{Rojo}}: $> 70$ dBA (nivel alto)
    \item \textcolor{blue}{\textbf{Azul}}: Modo calibraci\'on / warmup
\end{itemize}

\subsection{Disposici\'on en el aula}

Se colocan 6 micr\'ofonos ATOM Echo en una cuadr\'icula 2$\times$3 para cubrir la superficie del aula:

\begin{figure}[H]
\centering
\begin{tikzpicture}[
    mic/.style={circle, draw, fill=dbgreen!30, minimum size=1.2cm, font=\small\bfseries},
]
    \node[font=\small, fill=gray!20, minimum width=8cm, minimum height=0.6cm] at (3.5, 4.5) {PIZARRA};

    \node[mic] at (1, 3) {01};
    \node[mic] at (1, 1) {03};
    \node[mic] at (1, -1) {05};

    \node[mic] at (6, 3) {02};
    \node[mic] at (6, 1) {04};
    \node[mic] at (6, -1) {06};

    \draw[dashed, gray] (-0.5, -2) rectangle (7.5, 4);

    \node[font=\small, fill=gray!20, minimum width=0.6cm, minimum height=3cm, rotate=90] at (3.5, -2.5) {PUERTA};
\end{tikzpicture}
\caption{Disposici\'on de los 6 micr\'ofonos distribuidos en el aula}
\end{figure}

\section{M5Stack Core2 v1.1 (Micr\'ofono central)}

El Core2 se utiliza como micr\'ofono central opcional, ofreciendo informaci\'on visual en su pantalla LCD.

\begin{table}[H]
\centering
\begin{tabularx}{\textwidth}{lX}
\toprule
\textbf{Caracter\'istica} & \textbf{Especificaci\'on} \\
\midrule
MCU & ESP32 (240 MHz dual-core, 8 MB PSRAM) \\
Micr\'ofono & SPM1423 (PDM) \\
Pantalla & LCD IPS 320$\times$240 (16-bit color) \\
Identificador & \texttt{mic\_central} \\
Interfaz & Muestra dBA, RSSI WiFi, uptime, estado \\
\bottomrule
\end{tabularx}
\caption{Especificaciones del M5Stack Core2}
\end{table}

% ============================================================
% CAPITULO 3: Firmware
% ============================================================
\chapter{Firmware}

\section{Visi\'on general}

El firmware se ejecuta en los microcontroladores ESP32 y realiza las siguientes funciones:

\begin{enumerate}
    \item Captura de audio PCM a 16 kHz mediante I2S (M5Unified).
    \item Filtrado A-weighting IIR para ponderaci\'on frecuencial.
    \item C\'alculo de nivel RMS y conversi\'on a dBA.
    \item Acumulaci\'on de mediciones durante 5 segundos.
    \item Publicaci\'on MQTT con la media y el pico del intervalo.
    \item Indicaci\'on visual mediante LED RGB.
\end{enumerate}

\section{Filtro A-weighting (dsp.h)}

\subsection{Fundamento te\'orico}

La ponderaci\'on A (A-weighting) es un filtro frecuencial definido en la norma \textbf{IEC~61672} que modela la respuesta del o\'ido humano. Aten\'ua las frecuencias bajas ($< 500$ Hz) y muy altas ($> 6$ kHz) donde el o\'ido es menos sensible.

La implementaci\'on utiliza una cascada de \textbf{3 secciones biquad IIR} en forma \textit{Direct Form II Transposed}:

\begin{equation}
H(z) = H_1(z) \cdot H_2(z) \cdot H_3(z) \cdot G
\end{equation}

Donde $G = 1.24358$ es la ganancia de normalizaci\'on a 0 dB en 1 kHz.

\subsection{Secciones del filtro}

\begin{table}[H]
\centering
\begin{tabular}{clcc}
\toprule
\textbf{Secci\'on} & \textbf{Tipo} & \textbf{Frecuencia corte} & \textbf{Orden} \\
\midrule
Biquad 1 & Paso alto & 20.6 Hz & 2\textsuperscript{o} ($Q = 0.5$) \\
Biquad 2 & Paso alto & 107.7 Hz & 1\textsuperscript{er} \\
Biquad 3 & Paso alto & 737.9 Hz & 1\textsuperscript{er} \\
\bottomrule
\end{tabular}
\caption{Secciones del filtro A-weighting}
\end{table}

\textbf{Nota:} El polo doble a 12194 Hz se omite por estar por encima de la frecuencia de Nyquist (8 kHz a $f_s = 16$ kHz). Error t\'ipico: $< 2$ dB en el rango 100--8000 Hz.

\subsection{Implementaci\'on}

Cada secci\'on biquad procesa una muestra seg\'un:

\begin{equation}
y[n] = b_0 \cdot x[n] + s_1[n-1]
\end{equation}
\begin{equation}
s_1[n] = b_1 \cdot x[n] - a_1 \cdot y[n] + s_2[n-1]
\end{equation}
\begin{equation}
s_2[n] = b_2 \cdot x[n] - a_2 \cdot y[n]
\end{equation}

\begin{lstlisting}[language=C, caption={Procesado biquad (dsp.h)}]
static inline float biquad_process(Biquad* f, float x) {
    float y = f->b0 * x + f->s1;
    f->s1 = f->b1 * x - f->a1 * y + f->s2;
    f->s2 = f->b2 * x - f->a2 * y;
    return y;
}
\end{lstlisting}

\subsection{Coeficientes}

\begin{lstlisting}[language=C, caption={Coeficientes del filtro A-weighting para $f_s = 16000$ Hz}]
// Biquad 1: 2nd-order HP, fc=20.6Hz, Q=0.5
{ b0=0.991965, b1=-1.983931, b2=0.991965,
  a1=-1.983906, a2=0.983946 }

// Biquad 2: 1st-order HP, fc=107.7Hz
{ b0=0.979286, b1=-0.979286, b2=0.0,
  a1=-0.958573, a2=0.0 }

// Biquad 3: 1st-order HP, fc=737.9Hz
{ b0=0.873487, b1=-0.873487, b2=0.0,
  a1=-0.746974, a2=0.0 }
\end{lstlisting}

\section{Pipeline de medici\'on}

El proceso completo desde la captura de audio hasta el env\'io MQTT es:

\begin{enumerate}
    \item \textbf{Captura}: 1024 muestras PCM 16-bit a 16 kHz ($\approx 64$ ms por buffer).
    \item \textbf{Filtrado}: Cada muestra pasa por las 3 secciones biquad + ganancia.
    \item \textbf{RMS}: C\'alculo del valor eficaz del buffer filtrado.

    \begin{equation}
    \text{RMS} = \sqrt{\frac{1}{N} \sum_{i=0}^{N-1} x_{\text{filtrado}}[i]^2}
    \end{equation}

    \item \textbf{Conversi\'on a dB}:

    \begin{equation}
    \text{dBA} = 20 \cdot \log_{10}(\text{RMS}) + \text{DB\_OFFSET}
    \end{equation}

    Con l\'imites $[20, 120]$ dBA.

    \item \textbf{Acumulaci\'on}: Se acumulan $\sim$80 buffers durante 5 segundos.
    \item \textbf{Env\'io MQTT}: Media aritm\'etica del intervalo + pico m\'aximo.
\end{enumerate}

\section{Calibraci\'on}

\subsection{Offset est\'atico (DB\_OFFSET)}

Cada micr\'ofono tiene un par\'ametro \texttt{DB\_OFFSET} en \texttt{config.h} que permite compensar diferencias de sensibilidad:

\begin{lstlisting}[language=C]
#define DB_OFFSET 0.0  // Ajustar segun microfono
\end{lstlisting}

\subsection{Modo calibraci\'on interactivo}

Manteniendo pulsado el \textbf{Bot\'on A} durante 2 segundos se activa el modo calibraci\'on:

\begin{enumerate}
    \item El dispositivo mide durante 10 segundos.
    \item Calcula la media de todas las lecturas.
    \item Muestra por puerto serie el resultado y la f\'ormula de ajuste.
\end{enumerate}

\begin{equation}
\text{DB\_OFFSET}_{\text{nuevo}} = \text{DB\_OFFSET}_{\text{actual}} + (\text{dB}_{\text{referencia}} - \text{dB}_{\text{medido}})
\end{equation}

\subsection{Calibraci\'on por micro en el backend}

Adem\'as del offset en firmware, el backend dispone de una tabla de calibraci\'on que aplica correcciones sin necesidad de reflashear:

\begin{lstlisting}[language=Java, caption={Calibraci\'on en el backend (index.js)}]
const MIC_CALIBRATION = {
    "mic_central": 7.0,  // Resta 7 dB al central
};

function applyCalibration(mic, dbValue) {
    const offset = MIC_CALIBRATION[mic] || 0;
    return Math.round((dbValue - offset) * 10) / 10;
}
\end{lstlisting}

\section{Conectividad WiFi}

El firmware soporta dos modos de autenticaci\'on:

\begin{table}[H]
\centering
\begin{tabularx}{\textwidth}{lXX}
\toprule
\textbf{Par\'ametro} & \textbf{WPA2-Personal} & \textbf{WPA2-Enterprise} \\
\midrule
\texttt{WIFI\_ENTERPRISE} & \texttt{false} & \texttt{true} \\
\texttt{WIFI\_SSID} & Nombre de la red & Nombre de la red \\
\texttt{WIFI\_USER} & (vac\'io) & Usuario EAP-PEAP \\
\texttt{WIFI\_PASSWORD} & Clave WiFi & Clave EAP \\
\bottomrule
\end{tabularx}
\caption{Modos de autenticaci\'on WiFi}
\end{table}

Estrategia de reconexion: 90 intentos con 1 segundo entre cada uno. Cada 30 intentos realiza un ciclo completo de desconexi\'on y reconexion. Si agota los intentos, reinicia el ESP32.

\section{Protocolo MQTT}

\subsection{Topic}

\begin{verbatim}
aulas/{ROOM_ID}/{MIC_ID}/noise
\end{verbatim}

Ejemplo: \texttt{aulas/aula\_01/mic\_03/noise}

\subsection{Payload JSON}

\begin{lstlisting}[style=json, caption={Mensaje MQTT de un micr\'ofono}]
{
    "room": "aula_01",
    "mic": "mic_03",
    "db": 52.3,
    "peak": 67.1,
    "timestamp": 1773942550
}
\end{lstlisting}

\section{Configuraci\'on (config.h)}

\begin{lstlisting}[language=C, caption={Archivo de configuraci\'on completo}]
// WiFi
#define WIFI_ENTERPRISE false
#define WIFI_SSID "IOT2"
#define WIFI_USER ""
#define WIFI_PASSWORD "IOT_Enlaces_205"

// MQTT
#define MQTT_BROKER "broker.hivemq.com"
#define MQTT_PORT 1883

// Identificacion
#define ROOM_ID "aula_01"
#define MIC_ID "mic_06"

// Medicion
#define SEND_INTERVAL_MS 5000
#define SAMPLE_RATE 16000
#define SAMPLES_PER_READ 1024
#define DB_OFFSET 0.0

// Umbrales
#define THRESHOLD_LOW 50.0
#define THRESHOLD_HIGH 70.0
\end{lstlisting}

\section{Compilaci\'on (platformio.ini)}

\begin{lstlisting}[caption={Configuraci\'on PlatformIO para ATOM Echo S3R}]
[env:atom-echo-s3r]
platform = espressif32
board = esp32-s3-devkitc-1
framework = arduino
monitor_speed = 115200
build_flags =
    -DARDUINO_USB_CDC_ON_BOOT=1
    -DBOARD_HAS_PSRAM
    -DCORE_DEBUG_LEVEL=1
lib_deps =
    m5stack/M5Unified@^0.2.2
    knolleary/PubSubClient@^2.8
    bblanchon/ArduinoJson@^7.2.1
\end{lstlisting}

% ============================================================
% CAPITULO 4: Herramienta de flasheo
% ============================================================
\chapter{Herramienta de flasheo}

\section{flash.py}

Script Python para flashear m\'ultiples dispositivos ATOM Echo de forma automatizada.

\subsection{Caracter\'isticas}

\begin{itemize}
    \item Autodetecci\'on de puertos COM (CP210x, CH340, JTAG, USB gen\'erico).
    \item Modificaci\'on din\'amica de \texttt{config.h} antes de compilar.
    \item Soporte para flasheo secuencial de los 6 micr\'ofonos (\texttt{--all}).
    \item Credenciales WPA2-Enterprise prealmacenadas por micro.
    \item Configuraci\'on WiFi por l\'inea de comandos.
\end{itemize}

\subsection{Uso}

\begin{lstlisting}[language=bash, caption={Ejemplos de uso de flash.py}]
# Flashear un micro con credenciales Enterprise prealmacenadas
python flash.py mic_01

# Flashear con WiFi personalizado (WPA2-Personal)
python flash.py mic_01 --wifi MiRed MiClave

# Flashear con WiFi Enterprise (3 parametros)
python flash.py mic_01 --wifi ENLACES-A204 dam251v Shb*XEdv

# Flashear todos los micros
python flash.py --all --wifi MiRed MiClave

# Listar puertos COM disponibles
python flash.py --list
\end{lstlisting}

\subsection{Credenciales WPA2-Enterprise}

\begin{table}[H]
\centering
\begin{tabular}{lll}
\toprule
\textbf{MIC\_ID} & \textbf{Usuario} & \textbf{Contrase\~na} \\
\midrule
mic\_01 & dam251v & Shb*XEdv \\
mic\_02 & dam252v & blG1K1vT \\
mic\_03 & dam253v & 1B6\%CHbl \\
mic\_04 & dam254v & 8tE\%Phs@ \\
mic\_05 & dam255v & \%@CXFvVm \\
mic\_06 & dam256v & JTg71@*b \\
\bottomrule
\end{tabular}
\caption{Credenciales Enterprise por micr\'ofono}
\end{table}

% ============================================================
% CAPITULO 5: Backend
% ============================================================
\chapter{Backend}

\section{Arquitectura}

El backend es un servidor \textbf{Node.js} que integra:

\begin{itemize}
    \item \textbf{Express.js}: API REST (puerto 3001).
    \item \textbf{WebSocket} (ws): Actualizaciones en tiempo real (puerto 3002).
    \item \textbf{Cliente MQTT}: Recibe datos de los micr\'ofonos.
    \item \textbf{SQLite}: Almacenamiento persistente con modo WAL.
    \item \textbf{Broker MQTT integrado} (Aedes, opcional): Elimina la dependencia de Docker/Mosquitto.
\end{itemize}

\section{Base de datos}

\subsection{Esquema}

\begin{lstlisting}[language=SQL, caption={Esquema de la base de datos SQLite}]
CREATE TABLE noise_readings (
    id         INTEGER PRIMARY KEY AUTOINCREMENT,
    room       TEXT    NOT NULL,
    mic        TEXT    NOT NULL DEFAULT 'mic_01',
    db_level   REAL    NOT NULL,
    peak_level REAL    NOT NULL,
    timestamp  INTEGER NOT NULL,
    created_at TEXT    DEFAULT (datetime('now'))
);

CREATE INDEX idx_room_mic_timestamp
    ON noise_readings (room, mic, timestamp DESC);
\end{lstlisting}

\subsection{Optimizaciones}

\begin{itemize}
    \item \textbf{WAL mode}: Permite lecturas concurrentes durante escrituras.
    \item \textbf{Synchronous NORMAL}: Balance entre rendimiento y seguridad.
    \item \textbf{Prepared statements}: Consultas precompiladas para rendimiento.
    \item \textbf{\'Indice compuesto}: B\'usquedas r\'apidas por aula, micro y rango temporal.
\end{itemize}

\section{API REST}

\begin{table}[H]
\centering
\begin{tabularx}{\textwidth}{llX}
\toprule
\textbf{M\'etodo} & \textbf{Endpoint} & \textbf{Descripci\'on} \\
\midrule
GET & \texttt{/api/rooms} & Lista de aulas con resumen actual \\
GET & \texttt{/api/rooms/:id} & Detalle de un aula (todos sus micros) \\
GET & \texttt{/api/rooms/:id/history} & Hist\'orico de un aula (par\'ams: from, to) \\
GET & \texttt{/api/rooms/:id/mics/:micId/history} & Hist\'orico de un micro espec\'ifico \\
GET & \texttt{/api/stats} & Estad\'isticas globales (par\'am: since) \\
GET & \texttt{/api/health} & Estado del sistema (MQTT, uptime) \\
\bottomrule
\end{tabularx}
\caption{Endpoints de la API REST}
\end{table}

\subsection{Ejemplo de respuesta: /api/rooms}

\begin{lstlisting}[style=json, caption={Respuesta del endpoint /api/rooms}]
{
  "rooms": [
    {
      "room": "aula_01",
      "db": 43.9,
      "peak": 73.8,
      "micCount": 7,
      "onlineCount": 7,
      "mics": [
        {
          "room": "aula_01",
          "mic": "mic_central",
          "db": 38.1,
          "peak": 72.3,
          "timestamp": 1773942550,
          "online": true
        }
      ]
    }
  ]
}
\end{lstlisting}

\section{Agregaci\'on: mediana vs. media}

El valor de dB mostrado para cada aula se calcula como la \textbf{mediana} de los micr\'ofonos online, en lugar de la media aritm\'etica. Esto proporciona robustez frente a micr\'ofonos at\'ipicos (outliers):

\begin{lstlisting}[language=Java, caption={C\'alculo de mediana para agregaci\'on del aula}]
function median(values) {
    if (values.length === 0) return 0;
    const sorted = [...values].sort((a, b) => a - b);
    const mid = Math.floor(sorted.length / 2);
    return sorted.length % 2 !== 0
        ? sorted[mid]
        : (sorted[mid - 1] + sorted[mid]) / 2;
}

const medianDb = onlineMics.length > 0
    ? median(onlineMics.map(m => m.db))
    : 0;
\end{lstlisting}

\textbf{Justificaci\'on:} Con 7 micr\'ofonos, si uno de ellos (p.ej.\ mic\_central) lee sistem\'aticamente m\'as alto por estar m\'as cerca de la fuente de sonido, la mediana ignora este valor extremo y refleja mejor el nivel general del aula.

\section{Protocolo WebSocket}

La comunicaci\'on WebSocket permite actualizaciones en tiempo real sin polling:

\begin{table}[H]
\centering
\begin{tabularx}{\textwidth}{lX}
\toprule
\textbf{Tipo de mensaje} & \textbf{Descripci\'on} \\
\midrule
\texttt{init} & Estado completo de todas las aulas (al conectar) \\
\texttt{room\_update} & Actualizaci\'on de una aula (cada 5s o al cambiar estado) \\
\texttt{mic\_update} & Actualizaci\'on individual de un micr\'ofono \\
\bottomrule
\end{tabularx}
\caption{Tipos de mensaje WebSocket}
\end{table}

\section{Broker MQTT integrado}

Si se establece \texttt{USE\_INTERNAL\_BROKER=true}, el backend arranca un broker MQTT (Aedes) en el puerto 1883, eliminando la necesidad de Docker o Mosquitto externo:

\begin{lstlisting}[language=Java, caption={Broker MQTT integrado (broker.js)}]
const aedes = require("aedes")();
const net = require("net");
const PORT = parseInt(process.env.MQTT_INTERNAL_PORT || "1883");

const server = net.createServer(aedes.handle);
server.listen(PORT, () => {
    console.log(`[BROKER] MQTT broker en puerto ${PORT}`);
});
\end{lstlisting}

\section{Detecci\'on de micr\'ofonos offline}

Un micr\'ofono se considera \textbf{offline} si no env\'ia datos durante 15 segundos. El backend comprueba cada 5 segundos y notifica a los clientes WebSocket del cambio de estado.

\section{Variables de entorno}

\begin{lstlisting}[caption={Archivo .env.example}]
MQTT_BROKER=mqtt://localhost
MQTT_PORT=1883
HTTP_PORT=3001
WS_PORT=3002
DB_PATH=./data/noise.db
USE_INTERNAL_BROKER=true
\end{lstlisting}

\section{Dependencias}

\begin{table}[H]
\centering
\begin{tabular}{ll}
\toprule
\textbf{Paquete} & \textbf{Versi\'on} \\
\midrule
aedes & 0.50.0 \\
better-sqlite3 & 11.7.0 \\
express & 4.21.2 \\
mqtt & 5.10.4 \\
ws & 8.18.0 \\
dotenv & 16.4.7 \\
\bottomrule
\end{tabular}
\caption{Dependencias del backend}
\end{table}

% ============================================================
% CAPITULO 6: Frontend
% ============================================================
\chapter{Frontend}

\section{Arquitectura}

El frontend es una aplicaci\'on \textbf{Next.js 15} con el App Router de React 19. Utiliza:

\begin{itemize}
    \item \textbf{TypeScript}: Tipado est\'atico en todo el c\'odigo.
    \item \textbf{Tailwind CSS}: Estilos utility-first con colores personalizados.
    \item \textbf{Recharts}: Gr\'aficas interactivas de nivel de ruido.
    \item \textbf{WebSocket}: Actualizaciones en tiempo real.
    \item \textbf{REST fallback}: Polling cada 5s si WebSocket no est\'a disponible.
\end{itemize}

\section{Estructura de componentes}

\begin{table}[H]
\centering
\begin{tabularx}{\textwidth}{lX}
\toprule
\textbf{Componente} & \textbf{Funci\'on} \\
\midrule
\texttt{page.tsx} & Dashboard principal: grid de aulas, barra de estado \\
\texttt{layout.tsx} & Cabecera con logo, leyenda de colores, indicador de conexi\'on \\
\texttt{RoomCard.tsx} & Tarjeta por aula: valor dB grande, pico, sparkline, conteo de micros \\
\texttt{RoomDetailView.tsx} & Vista detallada: mapa 2$\times$3, micro central, estad\'isticas \\
\texttt{MicDetailView.tsx} & Detalle individual: valor dB, gr\'afica hist\'orica \\
\texttt{NoiseChart.tsx} & Gr\'afica Recharts: \'area dB + l\'inea pico + l\'ineas referencia \\
\texttt{useNoiseData.ts} & Hook personalizado: WebSocket, REST, navegaci\'on, estado \\
\bottomrule
\end{tabularx}
\caption{Componentes del frontend}
\end{table}

\section{Vistas del dashboard}

\subsection{Vista general (Dashboard)}

Muestra una cuadr\'icula responsive (1--4 columnas) con una tarjeta por aula. Cada tarjeta incluye:

\begin{itemize}
    \item N\'umero grande con el nivel actual en dBA (con animaci\'on pulse).
    \item Valor pico.
    \item Tiempo desde la \'ultima lectura (``hace X segundos'').
    \item Conteo de micr\'ofonos online/total.
    \item Sparkline con las \'ultimas 30 lecturas.
    \item Color de fondo seg\'un nivel (verde/amarillo/rojo).
\end{itemize}

Barra superior de estado: indicador de conexi\'on, total de aulas, media global, aulas por encima del umbral.

\subsection{Vista detallada del aula}

Al hacer clic en una tarjeta se accede a la vista detallada que muestra:

\begin{itemize}
    \item \textbf{Micr\'ofonos distribuidos}: Mapa 2$\times$3 con la posici\'on real de cada micro. C\'irculos coloreados con el valor dB. Los micros offline aparecen en gris.
    \item \textbf{Micr\'ofono central} (si existe \texttt{mic\_central}): C\'irculo grande con gradiente, comparativa entre media de los distribuidos y el central, diferencia.
    \item Tarjetas de estad\'isticas: media, pico m\'aximo, conteo de micros online.
    \item Navegaci\'on breadcrumb.
\end{itemize}

\subsection{Vista detallada del micr\'ofono}

Al hacer clic en un micr\'ofono individual:

\begin{itemize}
    \item Valor dB grande con fondo coloreado.
    \item Gr\'afica hist\'orica (Recharts) con \'area para dB, l\'inea para pico, y l\'ineas de referencia a 50 y 70 dB.
    \item Estad\'isticas: media, m\'aximo, m\'inimo del periodo.
\end{itemize}

\section{Hook useNoiseData}

El hook centraliza toda la gesti\'on de estado y comunicaci\'on:

\begin{lstlisting}[language=Java, caption={Interfaz del estado (useNoiseData.ts)}]
interface NoiseDataState {
    rooms: Map<string, RoomSummary>;
    connected: boolean;
    view: "dashboard" | "room" | "mic";
    selectedRoom: string | null;
    selectedMic: string | null;
    history: HistoryReading[];
    loading: boolean;
}
\end{lstlisting}

\textbf{Funciones expuestas:}
\begin{itemize}
    \item \texttt{selectRoom(id)}: Navega a la vista del aula.
    \item \texttt{selectMic(room, mic)}: Carga hist\'orico y navega a vista del micro.
    \item \texttt{goBack()}: Navega hacia atr\'as en la jerarqu\'ia.
    \item \texttt{getSparkline(room)}: Devuelve array de \'ultimos 30 valores dB.
\end{itemize}

\section{Estilos y colores}

\begin{table}[H]
\centering
\begin{tabular}{llll}
\toprule
\textbf{Nivel} & \textbf{Rango} & \textbf{Color} & \textbf{C\'odigo} \\
\midrule
Bajo & $< 50$ dBA & \textcolor{dbgreen}{Verde} & \texttt{\#22c55e} \\
Medio & $50-70$ dBA & \textcolor{dbyellow}{Amarillo} & \texttt{\#eab308} \\
Alto & $> 70$ dBA & \textcolor{dbred}{Rojo} & \texttt{\#ef4444} \\
\bottomrule
\end{tabular}
\caption{Escala de colores seg\'un nivel de ruido}
\end{table}

\section{Dependencias}

\begin{table}[H]
\centering
\begin{tabular}{ll}
\toprule
\textbf{Paquete} & \textbf{Versi\'on} \\
\midrule
next & 15.1.0 \\
react & 19.0.0 \\
recharts & 2.15.0 \\
tailwindcss & 3.4.16 \\
typescript & 5.7.2 \\
\bottomrule
\end{tabular}
\caption{Dependencias del frontend}
\end{table}

% ============================================================
% CAPITULO 7: Simulador
% ============================================================
\chapter{Simulador}

\section{Descripci\'on}

El simulador genera datos de ruido realistas sin necesidad de hardware f\'isico. \'Util para:

\begin{itemize}
    \item Demostraciones y presentaciones.
    \item Desarrollo y pruebas del frontend/backend.
    \item Validaci\'on de la cadena completa de datos.
\end{itemize}

\section{Escenarios realistas}

\begin{table}[H]
\centering
\begin{tabularx}{\textwidth}{lccX}
\toprule
\textbf{Escenario} & \textbf{dB base} & \textbf{$\pm$ dB} & \textbf{Descripci\'on} \\
\midrule
Clase tranquila & 35 & 8 & Explicaci\'on en voz baja \\
Explicaci\'on profesor & 50 & 10 & Clase magistral normal \\
Trabajo en grupo & 58 & 12 & Actividad colaborativa \\
Debate activo & 65 & 10 & Participaci\'on alta \\
Recreo & 72 & 15 & Hora de descanso \\
Examen & 28 & 5 & Silencio m\'aximo \\
Aula vac\'ia & 20 & 3 & Sin ocupantes \\
\bottomrule
\end{tabularx}
\caption{Escenarios de simulaci\'on}
\end{table}

\section{Uso}

\begin{lstlisting}[language=bash, caption={Ejecuci\'on del simulador}]
# Por defecto: 6 aulas, intervalo 2s
node simulate.js

# 12 aulas con intervalo de 1 segundo
node simulate.js --rooms 12 --interval 1000

# Broker personalizado
node simulate.js --broker mqtt://192.168.1.50
\end{lstlisting}

El simulador incluye 20 nombres de aulas predefinidos (aula\_101 a aula\_208, laboratorios, biblioteca, gimnasio, etc.) y genera transiciones suaves entre escenarios con picos aleatorios y decaimiento exponencial.

% ============================================================
% CAPITULO 8: Despliegue
% ============================================================
\chapter{Despliegue}

\section{Requisitos previos}

\begin{itemize}
    \item \textbf{Node.js} $\geq$ 18 (recomendado 20 LTS o 22).
    \item \textbf{Python} $\geq$ 3.8 con \texttt{platformio} y \texttt{pyserial} (solo para flasheo).
    \item \textbf{Docker} (opcional, solo si se usa Mosquitto externo).
\end{itemize}

\section{Escenario 1: Producci\'on con hardware real}

\subsection{Paso 1: Flashear los micr\'ofonos}

\begin{lstlisting}[language=bash]
cd firmware
pip install platformio pyserial
python flash.py mic_01  # Repetir para mic_02..06
\end{lstlisting}

\subsection{Paso 2: Arrancar el servidor}

\textbf{Windows:}
\begin{lstlisting}[language=bash]
start.bat
\end{lstlisting}

\textbf{Linux/Mac:}
\begin{lstlisting}[language=bash]
chmod +x start.sh && ./start.sh
\end{lstlisting}

Esto arranca:
\begin{enumerate}
    \item Backend con broker MQTT integrado (puertos 1883, 3001, 3002).
    \item Frontend en modo desarrollo (puerto 3000).
\end{enumerate}

\subsection{Paso 3: Acceder al dashboard}

\begin{itemize}
    \item Dashboard: \url{http://localhost:3000}
    \item API: \url{http://localhost:3001/api/rooms}
    \item MQTT: \texttt{localhost:1883}
\end{itemize}

\section{Escenario 2: Demo sin hardware}

\begin{lstlisting}[language=bash]
# Windows
scripts\demo.bat

# Linux/Mac
./scripts/demo.sh
\end{lstlisting}

Arranca el simulador en lugar de esperar datos de micr\'ofonos reales.

\section{Escenario 3: Docker}

\begin{lstlisting}[language=bash]
docker compose up -d    # Mosquitto + Backend
cd frontend && npm run dev  # Frontend aparte
\end{lstlisting}

\section{Puertos utilizados}

\begin{table}[H]
\centering
\begin{tabular}{lll}
\toprule
\textbf{Puerto} & \textbf{Servicio} & \textbf{Protocolo} \\
\midrule
1883 & Broker MQTT & TCP \\
3000 & Frontend (Next.js) & HTTP \\
3001 & Backend API (Express) & HTTP \\
3002 & WebSocket & WS \\
9001 & MQTT WebSocket (Mosquitto) & WS \\
\bottomrule
\end{tabular}
\caption{Puertos del sistema}
\end{table}

% ============================================================
% CAPITULO 9: Flujo de datos
% ============================================================
\chapter{Flujo de datos}

\section{Camino en tiempo real}

\begin{figure}[H]
\centering
\begin{tikzpicture}[
    node distance=1.2cm,
    box/.style={draw, rounded corners, minimum height=1cm, minimum width=3.5cm, align=center, fill=#1!15, font=\small},
    arrow/.style={-{Stealth[length=3mm]}, thick},
    label/.style={font=\scriptsize, midway, above},
]
    \node[box=blue] (mic) {ATOM Echo\\Captura + Filtro A};
    \node[box=orange, below=of mic] (mqtt) {Broker MQTT};
    \node[box=green, below=of mqtt] (backend) {Backend\\Calibraci\'on + SQLite};
    \node[box=green, right=3cm of backend] (ws) {WebSocket\\Server};
    \node[box=purple, below=of ws] (frontend) {Frontend\\Dashboard};

    \draw[arrow] (mic) -- node[right, font=\scriptsize] {WiFi / JSON} (mqtt);
    \draw[arrow] (mqtt) -- node[right, font=\scriptsize] {Subscribe} (backend);
    \draw[arrow] (backend) -- node[above, font=\scriptsize] {Broadcast} (ws);
    \draw[arrow] (ws) -- node[right, font=\scriptsize] {room\_update} (frontend);
\end{tikzpicture}
\caption{Flujo de datos en tiempo real}
\end{figure}

\section{Camino hist\'orico}

Cuando el usuario hace clic en un micr\'ofono para ver su hist\'orico:

\begin{enumerate}
    \item El frontend llama a \texttt{GET /api/rooms/:id/mics/:micId/history}.
    \item El backend consulta SQLite con el rango temporal.
    \item Devuelve array de lecturas con timestamps.
    \item El componente \texttt{NoiseChart} renderiza la gr\'afica.
\end{enumerate}

\section{Fallback por polling}

Si la conexi\'on WebSocket se pierde, el hook \texttt{useNoiseData} activa un polling REST cada 5 segundos a \texttt{/api/rooms} hasta que el WebSocket se reconecta (reintento autom\'atico cada 3 segundos).

% ============================================================
% CAPITULO 10: Estructura de ficheros
% ============================================================
\chapter{Estructura de ficheros}

\begin{lstlisting}[caption={\'Arbol completo del proyecto}, basicstyle=\ttfamily\scriptsize]
microfonoAula/
|-- firmware/                      # ATOM Echo S3R
|   |-- src/
|   |   |-- config.h              # Configuracion WiFi/MQTT/IDs
|   |   |-- dsp.h                 # Filtro A-weighting IIR
|   |   +-- main.cpp              # Logica principal
|   |-- platformio.ini            # Build config
|   +-- flash.py                  # Herramienta de flasheo
|
|-- firmware_core2/                # M5Stack Core2
|   |-- src/
|   |   |-- config.h
|   |   |-- dsp.h
|   |   +-- main.cpp              # Con LCD display
|   |-- platformio.ini
|   +-- flash_core2.py
|
|-- backend/                       # Servidor Node.js
|   |-- src/
|   |   |-- index.js              # API + WS + MQTT + SQLite
|   |   +-- broker.js             # Broker MQTT integrado
|   |-- package.json
|   |-- .env.example
|   +-- Dockerfile
|
|-- frontend/                      # Dashboard Next.js
|   |-- src/
|   |   |-- app/
|   |   |   |-- page.tsx          # Dashboard principal
|   |   |   |-- layout.tsx        # Cabecera
|   |   |   +-- globals.css       # Estilos globales
|   |   |-- components/
|   |   |   |-- RoomCard.tsx
|   |   |   |-- RoomDetailView.tsx
|   |   |   |-- MicDetailView.tsx
|   |   |   +-- NoiseChart.tsx
|   |   +-- lib/
|   |       +-- useNoiseData.ts   # Hook principal
|   |-- package.json
|   |-- next.config.ts
|   +-- tailwind.config.ts
|
|-- simulator/                     # Simulador sin hardware
|   |-- simulate.js
|   +-- package.json
|
|-- scripts/                       # Scripts de automatizacion
|   |-- demo.bat / demo.sh
|   |-- run_backend.bat
|   +-- run_frontend.bat
|
|-- start.bat / start.sh          # Arranque produccion
|-- setup.bat                      # Instalacion
|-- docker-compose.yml            # Docker (Mosquitto)
|-- mosquitto.conf                # Config Mosquitto
|-- requirements.txt              # Python deps
|-- README.md                      # Guia rapida
+-- DESPLIEGUE.md                 # Guia de despliegue
\end{lstlisting}

% ============================================================
% CAPITULO 11: Resumen de decisiones tecnicas
% ============================================================
\chapter{Decisiones t\'ecnicas}

\section{Mediana vs. media aritm\'etica}

Se utiliza la \textbf{mediana} para agregar los valores de los micr\'ofonos de un aula en lugar de la media aritm\'etica. Raz\'on: el micr\'ofono central, por su posici\'on m\'as cercana al profesor, registra sistem\'aticamente valores superiores (4--10 dB seg\'un el nivel de actividad). La mediana es robusta frente a estos outliers posicionales.

\section{A-weighting vs. ponderaci\'on plana}

Se implementa ponderaci\'on A (dBA) porque:
\begin{itemize}
    \item Es el est\'andar internacional para medici\'on de ruido ambiental (IEC~61672).
    \item Refleja la percepci\'on real del o\'ido humano.
    \item Aten\'ua el ruido de baja frecuencia (ventilaci\'on, tr\'afico) que no molesta tanto.
\end{itemize}

\section{Media de 5 segundos}

En lugar de enviar lecturas instant\'aneas, se acumula durante 5 segundos:
\begin{itemize}
    \item Reduce el tr\'afico MQTT ($\sim$12 mensajes/minuto por micro vs.\ $\sim$900 instant\'aneos).
    \item Suaviza fluctuaciones r\'apidas que no son representativas.
    \item El pico se preserva por separado para no perder eventos puntuales.
\end{itemize}

\section{WebSocket + REST fallback}

Se usa WebSocket como canal primario por su baja latencia, con fallback autom\'atico a polling REST. Esto garantiza que el dashboard funcione incluso en redes que bloquean WebSocket.

\section{SQLite vs. bases de datos externas}

Se eligi\'o SQLite por:
\begin{itemize}
    \item Cero configuraci\'on (fichero local).
    \item Rendimiento excelente para un solo servidor.
    \item Modo WAL para lecturas concurrentes.
    \item Portabilidad (la base de datos es un \'unico fichero \texttt{.db}).
\end{itemize}

\section{Broker MQTT integrado}

Se incluye Aedes como broker opcional para eliminar la dependencia de Docker/Mosquitto en despliegues simples. El backend puede actuar como broker y cliente MQTT simult\'aneamente.

\section{Calibraci\'on dual}

Se implementan dos niveles de calibraci\'on:
\begin{enumerate}
    \item \textbf{Firmware} (\texttt{DB\_OFFSET}): Compensaci\'on est\'atica por hardware. Requiere reflashear.
    \item \textbf{Backend} (\texttt{MIC\_CALIBRATION}): Compensaci\'on din\'amica por posici\'on. Ajuste inmediato sin tocar el hardware.
\end{enumerate}

% ============================================================
% Apendice
% ============================================================
\appendix

\chapter{Referencia r\'apida de comandos}

\begin{table}[H]
\centering
\begin{tabularx}{\textwidth}{lX}
\toprule
\textbf{Comando} & \textbf{Descripci\'on} \\
\midrule
\texttt{start.bat} & Arrancar backend + frontend (Windows) \\
\texttt{./start.sh} & Arrancar backend + frontend (Linux/Mac) \\
\texttt{scripts/demo.bat} & Demo con simulador (Windows) \\
\texttt{python flash.py mic\_01} & Flashear un micr\'ofono \\
\texttt{python flash.py --all} & Flashear todos los micros \\
\texttt{python flash.py --list} & Listar puertos COM \\
\texttt{node simulate.js} & Iniciar simulador (6 aulas) \\
\texttt{docker compose up -d} & Arrancar con Docker \\
\texttt{npm run dev} & Frontend en modo desarrollo \\
\texttt{npm run build} & Compilar frontend para producci\'on \\
\bottomrule
\end{tabularx}
\caption{Referencia r\'apida de comandos}
\end{table}

\chapter{Glosario}

\begin{description}[style=nextline]
    \item[dBA] Decibelios con ponderaci\'on A. Unidad de medida de nivel sonoro que tiene en cuenta la sensibilidad del o\'ido humano a distintas frecuencias.
    \item[A-weighting] Curva de ponderaci\'on frecuencial definida en IEC~61672 que modela la respuesta del o\'ido humano.
    \item[Biquad] Filtro digital IIR de segundo orden. Bloque b\'asico para construir filtros m\'as complejos.
    \item[RMS] Root Mean Square (valor eficaz). Medida de la amplitud de una se\~nal.
    \item[MQTT] Message Queuing Telemetry Transport. Protocolo de mensajer\'ia ligero para IoT.
    \item[I2S] Inter-IC Sound. Protocolo de comunicaci\'on serie para audio digital.
    \item[PDM] Pulse Density Modulation. T\'ecnica de modulaci\'on usada por micr\'ofonos MEMS.
    \item[WAL] Write-Ahead Logging. Modo de SQLite que permite lecturas concurrentes durante escrituras.
    \item[SPL] Sound Pressure Level. Nivel de presi\'on sonora medido en decibelios.
    \item[MEMS] Micro-Electro-Mechanical Systems. Tecnolog\'ia de fabricaci\'on de micr\'ofonos miniaturizados.
    \item[EAP-PEAP] Protected Extensible Authentication Protocol. M\'etodo de autenticaci\'on para WPA2-Enterprise.
\end{description}

\end{document}
